\documentclass{article}
\usepackage{amsmath,amsfonts,graphicx}
\usepackage[a4paper, left=2cm, right=2cm, top=2cm, bottom=2cm]{geometry}
\usepackage{hyperref}
\usepackage{listings}
\usepackage{xcolor}
\setlength{\parskip}{2pt}
\setlength{\parindent}{0pt}

\definecolor{codegray}{gray}{0.95}
\lstset{
	backgroundcolor=\color{codegray},
	basicstyle=\ttfamily\small,
	frame=single,
	breaklines=true
}

\pagestyle{empty}


\begin{document}
	
	% Titleplage
	\pagestyle{empty}
	\begin{center}
		\LARGE
		\textbf{Evaluation: Potential CLL-Classifiers} \vspace{0.5cm}
		\hrule
	\end{center}
	
	
	\tableofcontents
	\section{Summary}
		\subsection*{Exp01}
			\begin{itemize}
				\item Computed the RMA for the concatenated GSE13159 and GSE16455 superset (before I split train/ test)
				\item This is clear data leakage, because I transform train and test sets into a similar space, therefore test validations are invalid
				\item This is not how real-world ML prediction work... You usually take parameters from a train set to sale the test data
				\item This approach might be appropriate for biological processing, since it inherently removes noise 
			\end{itemize}
		
		\subsection*{Exp02}
			\begin{itemize}
				\item Classical Standard Scaling considering train and test set separate
				\item Meaning that I compute the test scaling from estimated parameters of the train set
			\end{itemize}
		
		\subsection*{Exp03}
			\begin{itemize}
				\item Classical Robust Scaling considering train and test set separate
				\item Meaning that I compute the test scaling from estimated parameters of the train set
			\end{itemize}
		
		\subsection*{Exp04}
			\begin{itemize}
				\item I got idea to approximate the RMA test sets with a shift-scale-vector estimated from train data
			\end{itemize}
			\begin{align*}
				\mu^{*} &= \mu^{[1]} - \mu^{[0]} \\
				\sigma^{*} &= \frac{\sigma^{[1]}}{\sigma^{[0]} + \epsilon} \\
				\implies& X^{scaled} = (X^{unscaled} + \mu^{*}) * \sigma^{*} \tag*{Approximation}
			\end{align*}
	
		\subsection{Conclusion}
			
			\begin{itemize}
				\item The most valid models seems to be a class weighted linear SVM model onto the approximated scaling
				\item I aim that this model didn't learn experimental artifact but rather true biological signals, further validated cross-experimentally
			\end{itemize}

			
	
	
	
	
	\newpage
	
	\section{Exp 01: Full Batch RMA}
	
		Since I introduced here intentionally strong data leakage, I will not investigate further after inspecting ROC curves of the trained model onto GSE131599-test.  \\
	
		\textbf{GSE131599-test (ROC)} 
		\begin{center}
			\begin{tabular}{ccc}
				KNN & Logit & Logit-weighted \\
				\includegraphics[scale=0.45]{"./images/01_bg_corrected_all_samples/roc_knn.png"} &
				\includegraphics[scale=0.45]{"./images/01_bg_corrected_all_samples/roc_logit.png"} &
				\includegraphics[scale=0.45]{"./images/01_bg_corrected_all_samples/roc_logit-weighted.png"} \\
				NB & svm-linear & svm-linear-weighted \\
				\includegraphics[scale=0.45]{"./images/01_bg_corrected_all_samples/roc_nb.png"} &
				\includegraphics[scale=0.45]{"./images/01_bg_corrected_all_samples/roc_svm-linear.png"} &
				\includegraphics[scale=0.45]{"./images/01_bg_corrected_all_samples/roc_svm-linear-weighted.png"} \\
				svm-rbf & svm-lar-weighted & \\
				\includegraphics[scale=0.45]{"./images/01_bg_corrected_all_samples/roc_nb.png"} &
				\includegraphics[scale=0.45]{"./images/01_bg_corrected_all_samples/roc_svm-rbf.png"} &
				\includegraphics[scale=0.45]{"./images/01_bg_corrected_all_samples/roc_svm-rbf-weighted.png"} \\
			\end{tabular}
		\end{center}
		\begin{itemize}
			\item I intentionally produce some data leakge, I wanted to see the effect for scaling experiments combined with RMA
			\item This approach is clearly not valid, since it produces over-promising results 
			\item Big Red Flag !
			\item Furthermore, also the prediction of GSE16455-test yields near perfect results, but i will not add those plots this report
			\item I will not further investigate here
		\end{itemize}
		\newpage
	
	
	
	\section{Exp 02: Standard Scaling}
	
	Most of the models ROC-AUC of the GSE131599-test data set was mostly around 0.5 and is therefore not a valid model for a diagnostic classifier. I only found that the Naive Bayes model reached a ROC-AUC of 0.293. I investigated further, since on could flip the prediction logic here, to get a ROC-AUC of rougly 0.7.\\
	
	\begin{center}
		\begin{tabular}{cc}
			GSE131599-test NB-ROC & GSE131599-test NB-PR \\
			\includegraphics[scale=0.5]{"./images/02_bg_corrected_mean_sd_scaled/roc_nb_test.png"} & 
			\includegraphics[scale=0.5]{"./images/02_bg_corrected_mean_sd_scaled/pr_nb_test.png"} \\
			GSE131599-train NB-ROC & GSE131599-train NB-PR \\
			\includegraphics[scale=0.5]{"./images/02_bg_corrected_mean_sd_scaled/roc_nb_train.png"} & 
			\includegraphics[scale=0.5]{"./images/02_bg_corrected_mean_sd_scaled/pr_nb_train.png"} \\
		\end{tabular}
	\end{center}
	\begin{itemize}
		\item Flipping the prediction logic, would be cherry-picking or the test data
		\item Also, it would make the train evaluation badly ... 
		\item Conclusion, Standard Scaling with a Naive Bayes doesn't yield a promising classifier
	\end{itemize}
	\newpage
	
	\section{Exp 03: Robust Scaling}
	I found that the SVM-rbf performed onto the GSE13159-test very reasonable, so I investigated its performance further. I included the insepction of the confusion matrix (CM), when using different optimal threshold evaluated onto the train.

	\begin{center}
		\begin{tabular}{cc}
			GSE131599-test SVM-rbf-ROC & GSE131599-test SVM-rbf-PR \\
			\includegraphics[scale=0.6]{"./images/03_bg_corrected_median_mad_scaled/roc_svm-rbf.png"} &
			\includegraphics[scale=0.6]{"./images/03_bg_corrected_median_mad_scaled/pr_svm-rbf.png"} \\
			GSE131599-test SVM-rbf-CM default & GSE131599-test SVM-rbf-CM gmean\\
			\includegraphics[scale=0.36]{"./images/03_bg_corrected_median_mad_scaled/cm_svm-rbf_d.png"} & 
			\includegraphics[scale=0.36]{"./images/03_bg_corrected_median_mad_scaled/cm_svm-rbf_g.png"} \\
			GSE131599-test SVM-rbf-CM Youden’s J statistic & GSE131599-test SVM-rbf-CM F1-Score \\
			\includegraphics[scale=0.36]{"./images/03_bg_corrected_median_mad_scaled/cm_svm-rbf_j.png"} & 
			\includegraphics[scale=0.36]{"./images/03_bg_corrected_median_mad_scaled/cm_svm-rbf_f1.png"} \\
		\end{tabular}
	\end{center}
	
	\begin{itemize}
		\item One can clearly reason, that the ROC and PRC are promising, but looking at the actual predictions, the model only outputs a single class. that's a big Red flag, since the model doesn't generalize well (even with default threshold)
	\end{itemize}
	
	\newpage
	
	
	\section{Exp 04: RMA-train with Approximation-test}
	
		I found that 3 models (Logit, SVM-lnear, SVM-linear-weighted) had an ROC-AUC $>0.6$. Of course, this might not be the best performance, but considering that I approximated the RPM for the test data with a shift-scale-vector, working on noisy biological, older technology (microarray) and without batch correction, I would argue that this is the right direction to go. Those models could be a classifiers that actually learned biological signals and not just experimental artifacts ... \\
	
		\textbf{- Logit}
		\begin{center}
			\begin{tabular}{cc}
				GSE13159-test Logit-ROC & GSE13159-test Logit-PR \\
				\includegraphics[scale=0.6]{"./images/04_bg_corrected_train_dqn_ext_test_scaled/roc_logit.png"} &
				\includegraphics[scale=0.6]{"./images/04_bg_corrected_train_dqn_ext_test_scaled/pr_logit.png"} \\
				GSE131599-test Logit-CM default & GSE131599-test Logit-CM gmean\\
				\includegraphics[scale=0.36]{"./images/04_bg_corrected_train_dqn_ext_test_scaled/cm_logit_d.png"} & 
				\includegraphics[scale=0.36]{"./images/04_bg_corrected_train_dqn_ext_test_scaled/cm_logit_g.png"} \\
				GSE131599-test Logit-CM Youden’s J statistic & GSE131599-test Logit-CM F1-Score \\
				\includegraphics[scale=0.36]{"./images/04_bg_corrected_train_dqn_ext_test_scaled/cm_logit_j.png"} & 
				\includegraphics[scale=0.36]{"./images/04_bg_corrected_train_dqn_ext_test_scaled/cm_logit_f1.png"} \\
			\end{tabular}
		\end{center}
		\begin{itemize}
			\item The CM reveals that the model predicts non-CLL class okish, but performance undesirably, predicting the CLL classes, what is for an diagnostic classifier not appropriate 
		\end{itemize}
		\newpage
		
		\textbf{- SVM linear}
		\begin{center}
			\begin{tabular}{cc}
				GSE13159-test SVM linear-ROC & GSE13159-test SVM linear-PR \\
				\includegraphics[scale=0.6]{"./images/04_bg_corrected_train_dqn_ext_test_scaled/roc_svm-linear.png"} &
				\includegraphics[scale=0.6]{"./images/04_bg_corrected_train_dqn_ext_test_scaled/pr_svm-linear.png"} \\
				GSE131599-test SVM linear-CM default & GSE131599-test SVM linear-CM gmean \\
				\includegraphics[scale=0.36]{"./images/04_bg_corrected_train_dqn_ext_test_scaled/cm_svm-linear_d.png"} & 
				\includegraphics[scale=0.36]{"./images/04_bg_corrected_train_dqn_ext_test_scaled/cm_svm-linear_g.png"} \\
				GSE131599-test SVM linear-CM Youden’s J statistic & GSE131599-testSVM linear-CM F1-Score \\
				\includegraphics[scale=0.36]{"./images/04_bg_corrected_train_dqn_ext_test_scaled/cm_svm-linear_j.png"} & 
				\includegraphics[scale=0.36]{"./images/04_bg_corrected_train_dqn_ext_test_scaled/cm_svm-linear_f1.png"} \\
			\end{tabular}
		\end{center}
		\begin{itemize}
			\item The default threshold yield a very good non-CLL classifier, but doesn't perform well on the CLL classfiication
			\item Two threshold yields a kind of good direction, with $57.1\%$ CLL correct and $64.9\%$ non-CLL correct
			\item But i would sanity check the train ROC and the CM for the both threshold next
		\end{itemize}
		
		\begin{center}
			\begin{tabular}{cc}
				GSE13159-train SVM linear-ROC & GSE13159-train SVM linear-PR \\
				\includegraphics[scale=0.6]{"./images/04_bg_corrected_train_dqn_ext_test_scaled/roc_svm-linear_train.png"} &
				\includegraphics[scale=0.6]{"./images/04_bg_corrected_train_dqn_ext_test_scaled/pr_svm-linear_train.png"} \\
				GSE131599-train SVM linear-CM default & GSE131599-train SVM linear-CM gmean \\
				\includegraphics[scale=0.36]{"./images/04_bg_corrected_train_dqn_ext_test_scaled/cm_svm-linear_d_train.png"} & 
				\includegraphics[scale=0.36]{"./images/04_bg_corrected_train_dqn_ext_test_scaled/cm_svm-linear_g_train.png"} \\
				GSE131599-train SVM linear-CM Youden’s J statistic & GSE131599-train SVM linear-CM F1-Score \\
				\includegraphics[scale=0.36]{"./images/04_bg_corrected_train_dqn_ext_test_scaled/cm_svm-linear_j_train.png"} & 
				\includegraphics[scale=0.36]{"./images/04_bg_corrected_train_dqn_ext_test_scaled/cm_svm-linear_f1_train.png"} \\
			\end{tabular}
		\end{center}
				
		\begin{itemize}
			\item The classifier performs perfect onto the train data (data seems to be separable with an hyperplane)
			\item Using the gmean threshold, yield a classifer that only predicts non-cll, so not valid at all
			\item But using the Youden's statistic threshold, yield a reasonable trainining classification, and would align with the performance of the GSE13159-test 
			\item So this model and this scaling seems to produce trustworthy predictions
			\item I would next inspect the GSE16355-test performance, 
		\end{itemize}
				
		\begin{center}
			\begin{tabular}{c}
				GSE16455-test SVM-linear Raw-Probabilities for samples \\
				\includegraphics[scale=0.7]{"./images/04_bg_corrected_train_dqn_ext_test_scaled/prob_svm-linear_f1_test2.png"}
			\end{tabular}
		\end{center}
		
		\begin{itemize}
			\item GSE16455-test with acc(default) = 0.59 is fine, but model has very low true positives on the GSE13159-test (not valid)
			\item GSE16455-test with acc(gmean) = 0.41, yield a okisch classification for GSE13159-test, but only predicts non-CLL onto GSE13159-train (not valid)
			\item GSE16455-test with acc(J stat.) = 0.41, yield a okisch classification for GSE13159-test, also learned pattern in train data (most valid)
			\item GSE16455-test with acc(f1) = 1.0, only outputs  non-CLL classification for GSE13159-test
		\end{itemize}
		
		\textbf{Conclusion}
		\begin{itemize}
			\item I would argue that the SVM-linear onto the approxiamted scaling truely learns biological signals from this inherently noise data 
			\item I can see that the shift-scaled approxiamtion of the RMA seems to work and performs better then just a standard or robust or standard scaling 
			\item Considering the GSE13159-train, GSE13159-test and GSE16455-test evaluation, it seems that the optimal threshold for the Yourdan's statistics performs best  
		\end{itemize}
		
		
		\newpage
		
		\textbf{- SVM linear weighted}
		\begin{center}
			\begin{tabular}{cc}
				GSE13159-test SVM linear-ROC & GSE13159-test SVM linear-PR \\
				\includegraphics[scale=0.6]{"./images/04_bg_corrected_train_dqn_ext_test_scaled/roc_svm-linear-weighted.png"} &
				\includegraphics[scale=0.6]{"./images/04_bg_corrected_train_dqn_ext_test_scaled/pr_svm-linear-weighted.png"} \\
				GSE13159-test SVM linear-CM default & GSE13159-test SVM linear-CM gmean \\
				\includegraphics[scale=0.36]{"./images/04_bg_corrected_train_dqn_ext_test_scaled/cm_svm-linear-weighted_d.png"} & 
				\includegraphics[scale=0.36]{"./images/04_bg_corrected_train_dqn_ext_test_scaled/cm_svm-linear-weighted_g.png"} \\
				GSE131599-test SVM linear-CM Youden’s J statistic & GSE131599-test SVM linear-CM F1-Score \\
				\includegraphics[scale=0.36]{"./images/04_bg_corrected_train_dqn_ext_test_scaled/cm_svm-linear-weighted_j.png"} & 
				\includegraphics[scale=0.36]{"./images/04_bg_corrected_train_dqn_ext_test_scaled/cm_svm-linear-weighted_f1.png"} \\
			\end{tabular}
		\end{center}
		\begin{itemize}
			\item The CM reveals that the model predicts reasonable with the f1-score optimizing threshold, the rest only output CLL classes 
			\item I noticed that the correct classification of CLL classes is close to 0.8 but only roughly 0.4 non-CLL classes were labeled correctly
			\item Considering these numbers, this would be a better medical classifier, then the SVM-linear without class weights, because detecting CLL classes correctly, might be more relevant (diagnostically speaking), then non-CLL classes
		\end{itemize}
	
		\newpage
	
		\begin{center}
			\begin{tabular}{cc}
				GSE13159-train SVM linear-ROC & GSE13159-train SVM linear-PR \\
				\includegraphics[scale=0.6]{"./images/04_bg_corrected_train_dqn_ext_test_scaled/roc_svm-linear-weighted_train.png"} &
				\includegraphics[scale=0.6]{"./images/04_bg_corrected_train_dqn_ext_test_scaled/pr_svm-linear-weighted_train.png"} \\
				GSE131599-train SVM linear-CM default & GSE131599-train SVM linear-CM gmean \\
				\includegraphics[scale=0.36]{"./images/04_bg_corrected_train_dqn_ext_test_scaled/cm_svm-linear-weighted_d_train.png"} & 
				\includegraphics[scale=0.36]{"./images/04_bg_corrected_train_dqn_ext_test_scaled/cm_svm-linear-weighted_g_train.png"} \\
				GSE131599-train SVM linear-CM Youden’s J statistic & GSE131599-train SVM linear-CM-weighted F1-Score \\
				\includegraphics[scale=0.36]{"./images/04_bg_corrected_train_dqn_ext_test_scaled/cm_svm-linear-weighted_j_train.png"} & 
				\includegraphics[scale=0.36]{"./images/04_bg_corrected_train_dqn_ext_test_scaled/cm_svm-linear-weighted_f1_train.png"} \\
			\end{tabular}
		\end{center}
		\begin{itemize}
			\item The CM reveals that the model predicts reasonable with the f1-score optimizing threshold for the same threshold then the GSE13159-test
		\end{itemize}
		
		\begin{center}
			\begin{tabular}{c}
				GSE16455-test SVM-linear-weighted Raw-Probabilities for samples \\
				\includegraphics[scale=0.7]{"./images/04_bg_corrected_train_dqn_ext_test_scaled/prob_svm-linear-weighted_f1_test2.png"}
			\end{tabular}
		\end{center}
		\begin{itemize}
			\item GSE16455-test with acc(default) = 0.0 (not valid)
			\item GSE16455-test with acc(gmean) = 0.0, (not valid)
			\item GSE16455-test with acc(J stat.) = 0.0 (not valid)
			\item GSE16455-test with acc(f1) = 0.53, this seems to be the most valid and promising model for this task,  $CM_{train}(\text{CLL}) = 0.987; CM_{train}(\text{non-CLL}) = 0.61; CM_{test}(\text{CLL}) = 0.786; CM_{test}(\text{non-CLL}) = 0.404$
		\end{itemize}
		
		\textbf{Conclusion}
		\begin{itemize}
			\item The SVM-linear with class weights seem to be the most valid model
			\item Especially the fact, that for an medical classifier, the TP are high, without cherry picking for threshold, just using an mathematical optimal threshold
		\end{itemize}
		
\end{document}
